Düşmeler; yaşlı bireylerde yaralanmanın, bağımsızlık kaybının ve ölümün başlıca
nedenlerinden biridir. Dünya genelinde yaşlı nüfusun hızla artması, düşme olaylarının
zamanında ve güvenilir biçimde tespit edilmesini bir kamu sağlığı önceliği hâline
getirmektedir. Giyilebilir ataletsel sensörler (ivmeölçer ve jiroskop), bireyin günlük
yaşamını kısıtlamadan sürekli izleme imkânı sunduğundan, otomatik düşme tespiti için
pratik bir donanım altyapısı sağlar. Bu problemde tespit edilemeyen bir düşmenin
(yanlış negatif) maliyeti, yanlış bir alarmın (yanlış pozitif) maliyetinden çok daha
yüksektir; çünkü kaçırılan bir düşme, yerde uzun süre yardımsız kalma ve buna bağlı
ciddi komplikasyonlar anlamına gelebilir. Dolayısıyla düşme tespiti doğası gereği geri
çağırma (recall) odaklı bir sınıflandırma problemidir.

Bu tez, gecikme gömme (delay embedding) tabanlı bir işlem hattı aracılığıyla kısa
giyilebilir ivmeölçer pencerelerine uygulanan kalıcı homolojinin (persistent homology);
değerlendirmesi titiz, öznitelik inşası ilkeli ve uç (edge) cihazlarda çalışabilecek
kadar hafif bir geri-çağırma odaklı düşme tespitini destekleyip desteklemediğini
araştırmaktadır. Çalışmanın merkezindeki iddia, başarının büyük ölçüde sınıflandırıcı
seçiminden değil, sinyali topolojik bir öznitelik uzayına taşıyan temsil haritasından
kaynaklandığıdır.

\vspace{2mm}
\noindent\textbf{Yöntem.}
Önerilen işlem hattı, her üç-eksenli kaydı önce yönelimden bağımsız bir sinyal büyüklük
vektörüne (SMV) dönüştürür. Tüm kayıtlar ortak bir örnekleme frekansı olan 50 Hz'e
yeniden örneklenir ve sabit uzunlukta pencerelere bölünür. Her pencerenin faz uzayı,
Takens'in gömme teoremi anlamında gecikme gömmeyle yeniden inşa edilir; böylece tek
boyutlu skaler sinyalden çok boyutlu bir nokta bulutu elde edilir. Bu nokta bulutu
üzerine bir filtrasyon kurulur ve sıfırıncı ile birinci boyutta kalıcı homoloji
hesaplanır. Elde edilen kalıcılık diyagramlarının (persistence diagram) her biri on
istatistikle özetlenir; bu istatistikler dört sinyal düzeyi tanımlayıcısıyla birleşerek
yirmi dört boyutlu, kompakt bir öznitelik vektörü ($\phi_\lambda$) oluşturur. Bu vektör,
kasıtlı olarak basit tutulan iki sınıflandırıcıya verilir: lojistik regresyon (LR) ve
destek vektör makinesi (SVM). Her ikisi de sınıf dengesizliğini gözetmek üzere dengeli
sınıf ağırlıklarıyla, SVM ayrıca Platt olasılık kalibrasyonuyla eğitilir.

Topolojik temsilin davranışı; pencere uzunluğu, gömme boyutu, gecikme, filtrasyon türü
(Vietoris--Rips, Alpha, SparseRips), uzaklık metriği ve eşik gibi hiperparametrelerden
oluşan bir demete ($\lambda$) bağlıdır. Bu yüksek boyutlu ve koşullu arama uzayı,
ızgara aramasını hesaplama açısından olanaksız kıldığından, hiperparametreler ağaç
yapılı Parzen kestiricisi (TPE) ile yürütülen Bayesçi optimizasyon kullanılarak
seçilir. Her veri kümesi için ayrı bir arama yürütülür ve sınıflandırıcıdan bağımsız
bir amaç fonksiyonu ile en iyi yapılandırma belirlenir.

\vspace{2mm}
\noindent\textbf{Veri Kümeleri ve Değerlendirme.}
Yöntem, üç açık karşılaştırma veri kümesi üzerinde değerlendirilmiştir: MobiFall,
SisFall ve FallAllD veri kümesinin 40 Hz türetilmiş çeşidi olan FAD\_40Hz. Bu veri
kümeleri; sensör konumu, örnekleme frekansı, denek profili ve düşme türleri bakımından
farklılaşarak yöntemin genellenebilirliğini sınamaya olanak tanır. Değerlendirme,
denek-duyarlı bir çerçevede üç protokolle yapılır: rastgele bölünmeye dayanan ve üst
sınır referansı sağlayan naif protokol; her deneğin tamamen dışarıda bırakıldığı,
birincil protokol olan denek-dışı-bırakma (leave-one-subject-out, LOSO) protokolü; ve
her deneğin yalnızca kendi verisiyle modellendiği kişisel protokol. Bu protokolleri,
karar eşiğinin geniş bir aralıkta taranmasına dayanan bir eşik analizi tamamlar. Birincil
başarım ölçütü, geri çağırmayı kesinliğe göre dört kat ağırlıklandıran $F_2$ skorudur;
denekler arası dağılım, $t$-dağılımına dayalı yüzde 95 güven aralıklarıyla raporlanır.
Tüm ölçekleme ve eşik parametreleri yalnızca eğitim katmanından kestirilerek bilgi
sızıntısı önlenir.

\vspace{2mm}
\noindent\textbf{Bulgular.}
LOSO protokolü altında temsil; SisFall'da $0.97$'ye, FAD\_40Hz'de $0.96$'ya ve
MobiFall'da $0.90$'a varan ortalama $F_2$ skoru elde etmiş, geri çağırma değeri ise tüm
veri kümesi--sınıflandırıcı bileşimlerinde $0.95$'in üzerinde kalmıştır. Bu sonuçlar,
temsilin önceden belirlenmiş başarım hedeflerini en az bir sınıflandırıcı için
karşıladığını göstermektedir. Bu etkinliğin kuramsal dayanağı Cohen--Steiner kararlılık
teoremidir; bu teorem, iki kalıcılık diyagramı arasındaki darboğaz (bottleneck)
uzaklığının, girdi sinyallerindeki bozulmanın supremum normuyla sınırlı olduğunu garanti
eder ve dolayısıyla elle tasarlanmış istatistiksel, spektral veya derin öğrenme tabanlı
öznitelik ailelerinde benzeri bulunmayan biçimsel bir gürültü-dayanıklılığı sertifikası
sağlar.

Üç veri kümesi yapısal olarak birbirinden farklı en iyi yapılandırmalara yakınsamıştır:
MobiFall geniş bir pencerede SparseRips kompleksini, SisFall kısa bir pencerede Alpha
kompleksini, FAD\_40Hz ise kısa bir pencerede standart Rips kompleksini seçmiştir. Bu
durum, evrensel bir yapılandırmanın var olmadığını ve en iyi topolojik geometrinin
düşme dinamiklerindeki ve nokta-bulutu yapısındaki gerçek farkları yansıttığını ortaya
koyar. Ayrıca iki sınıflandırıcı arasındaki $F_2$ farkı tüm veri kümelerinde $0.02$'nin
altında kalırken, topolojik hiperparametrelerin ayarlanması $F_2$'yi büyüklük olarak çok
daha geniş bir aralıkta değiştirmektedir. Bu gözlem, başarımın sınıflandırıcı sınıfından
çok temsil haritası tarafından belirlendiği temel iddiasını güçlü biçimde
desteklemektedir.

Son olarak, model yeniden eğitimi gerektirmeyen sonradan (post-hoc) eşik
kişiselleştirmesinin etkisi incelenmiştir. Denek düzeyindeki kazanım dağılımı, tüm veri
kümesi--sınıflandırıcı bileşimlerinde sağa çarpıktır: deneklerin çoğunluğu yalnızca
marjinal fayda sağlarken küçük bir azınlık $F_2$'de $0.10$'dan büyük kazanımlar elde
etmektedir. Bu heterojenlik, tek bir sabit eşik kaymasının tüm denekleri eşit biçimde
iyileştireceği varsayımını çürütmekte ve en uygun eşiğin denek-özgü olduğunu
göstermektedir.

\vspace{2mm}
\noindent\textbf{Sonuç.}
Tezin başlıca katkıları şunlardır: (i) giyilebilir sinyallerden geri-çağırma odaklı
düşme tespiti için kuramsal olarak güvenceli, hesaplama açısından hafif bir topolojik
temsil haritası; (ii) bu temsili üç ayrı veri kümesinde tutarlı biçimde optimize eden
Bayesçi arama yordamı; (iii) naif, LOSO ve kişisel protokolleri bir araya getiren
denek-duyarlı bir değerlendirme çerçevesi; ve (iv) model yeniden eğitimi gerektirmeyen,
çıkarım zamanında uygulanabilen ve denek-özgü heterojenliği ortaya koyan bir eşik
kişiselleştirme analizi. Bulgular bütün olarak, topolojik temsilin basit
sınıflandırıcılarla dahi güçlü ve gürültüye dayanıklı düşme tespiti sağlayabileceğini;
ancak kişiselleştirmenin homojen bir iyileştirme sunmadığını ve denek bazında ele
alınması gerektiğini göstermektedir.

\vspace{3mm}\noindent
\textbf{Anahtar Kelimeler:} düşme tespiti, giyilebilir sensörler, topolojik veri analizi,
kalıcı homoloji, gecikme gömme, geri-çağırma odaklı sınıflandırma.
